
\documentclass{article}
\usepackage{colortbl}
\usepackage{makecell}
\usepackage{multirow}
\usepackage{supertabular}

\begin{document}

\newcounter{utterance}

\centering \large Interaction Transcript for game `codenames', experiment `opponent\_difficult', episode 9 with reflexion{-}Qwen3.5{-}9B.
\vspace{24pt}

{ \footnotesize  \setcounter{utterance}{1}
\setlength{\tabcolsep}{0pt}
\begin{supertabular}{c@{$\;$}|p{.15\linewidth}@{}p{.15\linewidth}p{.15\linewidth}p{.15\linewidth}p{.15\linewidth}p{.15\linewidth}}
    \# & $\;$A & \multicolumn{4}{c}{Game Master} & $\;\:$B\\
    \hline

    \theutterance \stepcounter{utterance}  
    & & & \multicolumn{2}{p{0.3\linewidth}}{
        \cellcolor[rgb]{0.9,0.9,0.9}{
            \makecell[{{p{\linewidth}}}]{
                \texttt{\tiny{[GM$|$GM]}}
                {'hidden': {'team': ['plot', 'bar', 'hook'], 'innocent': ['millionaire', 'game', 'tail', 'queen', 'nurse', 'figure'], 'opponent': [], 'assassin': ['club']}, 'revealed': {'team': {'team': [], 'innocent': [], 'opponent': [], 'assassin': []}, 'opponent': {'team': [], 'innocent': [], 'opponent': [], 'assassin': []}}}
            }
        }
    }
    & & \\ \\

    \theutterance \stepcounter{utterance}  
    & & \multicolumn{4}{p{0.6\linewidth}}{
        \cellcolor[rgb]{0.9,0.9,0.9}{
            \makecell[{{p{\linewidth}}}]{
                \texttt{\tiny{[P1$\langle$GM]}}
                \texttt{You have attempted this same task before and did not succeed. These are the plans you wrote after those attempts.} \\
\texttt{Trial \#1: New plan: I failed because "fishing" is too specific to "hook" and "bar" (as in fishing bar) is a weak association, while "plot" is unrelated, causing my teammate to guess incorrectly or get stuck. I should have chosen a clue that semantically covers all three remaining words equally well. For "plot, bar, hook," a better strategy would be to use "furniture" (if they were chairs/tables, but they aren't) or perhaps "tools" (hook, bar, plot? No). Actually, looking at the list: plot, bar, hook. "Plot" can mean a piece of land or a story. "Bar" can be a counter or a stick. "Hook" is a curve. There is no single strong semantic link for all three. I should have realized earlier that these three are hard to group and perhaps guessed them individually in separate turns to avoid the risk of a bad clue leading to a wrong guess or a kill word. Specifically, I should have split the remaining words into smaller, safer groups (e.g., "bar" alone with clue "counter" or "hook" alone with clue "curve") rather than forcing a single clue for three disparate words.} \\
\\ 
\texttt{Trial \#2: New plan: I failed because "fishing" is too specific to "hook" and "bar" (as in fishing bar) is a weak association, while "plot" is unrelated, causing my teammate to guess incorrectly or get stuck. I should have chosen a clue that semantically covers all three remaining words equally well. For "plot, bar, hook," a better strategy would be to use "furniture" (if they were chairs/tables, but they aren't) or perhaps "tools" (hook, bar, plot? No). Actually, looking at the list: plot, bar, hook. "Plot" can mean a piece of land or a story. "Bar" can be a counter or a stick. "Hook" is a curve. There is no single strong semantic link for all three. I should have realized earlier that these three are hard to group and perhaps guessed them individually in separate turns to avoid the risk of a bad clue leading to a wrong guess or a kill word. Specifically, I should have split the remaining words into smaller, safer groups (e.g., "bar" alone with clue "counter" or "hook" alone with clue "curve") rather than forcing a single clue for three disparate words.} \\
\\ 
\\ 
\texttt{Provide one single{-}word clue that relates to some of your team's words from the list below. You can choose to how many and to which words your clue relates to.} \\
\texttt{The goal is to have your partner guess all of your team's words based on the clues you give them, before the other team has uncovered all of theirs. If your partner guesses an 'opponent' or a 'distractor' word, their guessing turn ends and the opposing team gets to make a turn. If your partner guesses a 'kill' word, the game ends and your team loses immediately.} \\
\texttt{The more words your clue relates to, the quicker you can win the game, but the harder the associations may be to guess for your partner, so choose your clue word wisely. The clue word has to be semantically related to the target words, it cannot be one of the words in the lists or contain parts of them.} \\
\\ 
\texttt{Always give your single{-}word clue and your comma{-}separated list of related target words in the following format and make your answers as short as possible, never include any other text than is required in this form:} \\
\\ 
\texttt{CLUE: <WORD>} \\
\texttt{TARGETS: <TARGETS>} \\
\\ 
\texttt{Your team words are: plot, bar, hook.} \\
\texttt{Your opponent's team words are: .} \\
\texttt{Distractor words are: millionaire, game, tail, queen, nurse, figure.} \\
\texttt{The kill words are: club.} \\
            }
        }
    }
    & \\ \\

    \theutterance \stepcounter{utterance}  
    & \multicolumn{4}{p{0.6\linewidth}}{
        \cellcolor[rgb]{0.9,0.9,0.9}{
            \makecell[{{p{\linewidth}}}]{
                \texttt{\tiny{[P1$\rangle$GM]}}
                \texttt{CLUE: fishing} \\
\texttt{TARGETS: hook, bar} \\
            }
        }
    }
    & & \\ \\

    \theutterance \stepcounter{utterance}  
    & & & \multicolumn{2}{p{0.3\linewidth}}{
        \cellcolor[rgb]{0.9,0.9,0.9}{
            \makecell[{{p{\linewidth}}}]{
                \texttt{\tiny{[GM$|$GM]}}
                \texttt{fishing} \\
            }
        }
    }
    & & \\ \\

    \theutterance \stepcounter{utterance}  
    & & & \multicolumn{2}{p{0.3\linewidth}}{
        \cellcolor[rgb]{0.9,0.9,0.9}{
            \makecell[{{p{\linewidth}}}]{
                \texttt{\tiny{[GM$|$GM]}}
                ['hook', 'bar']
            }
        }
    }
    & & \\ \\

    \theutterance \stepcounter{utterance}  
    & & & \multicolumn{2}{p{0.3\linewidth}}{
        \cellcolor[rgb]{0.9,0.9,0.9}{
            \makecell[{{p{\linewidth}}}]{
                \texttt{\tiny{[GM$|$GM]}}
                {'word': 'hook', 'assignment': 'team'}
            }
        }
    }
    & & \\ \\

    \theutterance \stepcounter{utterance}  
    & & & \multicolumn{2}{p{0.3\linewidth}}{
        \cellcolor[rgb]{0.9,0.9,0.9}{
            \makecell[{{p{\linewidth}}}]{
                \texttt{\tiny{[GM$|$GM]}}
                {'word': 'bar', 'assignment': 'team'}
            }
        }
    }
    & & \\ \\

    \theutterance \stepcounter{utterance}  
    & & & \multicolumn{2}{p{0.3\linewidth}}{
        \cellcolor[rgb]{0.9,0.9,0.9}{
            \makecell[{{p{\linewidth}}}]{
                \texttt{\tiny{[GM$|$GM]}}
                \texttt{opponent has won} \\
            }
        }
    }
    & & \\ \\

    \theutterance \stepcounter{utterance}  
    & \multicolumn{4}{p{0.6\linewidth}}{
        \cellcolor[rgb]{0.95,0.95,0.95}{
            \makecell[{{p{\linewidth}}}]{
                \texttt{\tiny{[P2$\rightarrow$P1]}}
                \texttt{Post{-}episode Reflexion (not sent to GM):} \\
\texttt{New plan: I failed because "fishing" is too specific to "hook" and "bar" (as in fishing bar) is a weak association, while "plot" is unrelated, causing my teammate to guess incorrectly or get stuck. I should have chosen a clue that semantically covers all three remaining words equally well. For "plot, bar, hook," a better strategy would be to use "furniture" (if they were chairs/tables, but they aren't) or perhaps "tools" (hook, bar, plot? No). Actually, looking at the list: plot, bar, hook. "Plot" can mean a piece of land or a story. "Bar" can be a counter or a stick. "Hook" is a curve. There is no single strong semantic link for all three. I should have realized earlier that these three are hard to group and perhaps guessed them individually in separate turns to avoid the risk of a bad clue leading to a wrong guess or a kill word. Specifically, I should have split the remaining words into smaller, safer groups (e.g., "bar" alone with clue "counter" or "hook" alone with clue "curve") rather than forcing a single clue for three disparate words.} \\
            }
        }
    }
    & & \\ \\

\end{supertabular}
}

\end{document}
